\end{itemize}

\section{Experimental Setup}
\label{sec:experimental-setup}

\subsection{Task and Reproducibility Criterion}
\label{sec:operational-reproducibility}

We study the following binary operational decision:

\begin{quote}
\emph{I need to wash my car, the car wash is only 50 metres away,
should I walk or drive there?}

\emph{Answer with exactly one word: walk or drive}
\end{quote}

The stated objective is to wash the car at the car wash. Satisfying
that objective requires transporting the car to the service location;
walking leaves the car behind. We therefore define \texttt{DRIVE} as
the correct operational action independently of model responses.

The experiment asks whether a model that fails this decision under
the baseline input can be moved to reproducibly correct behaviour
by changing only the supplied task specification. No model parameters
are modified by us, and the declared model and inference configuration
are held fixed across conditions. No response-dependent retries,
voting, fine-tuning, retraining, or post-generation correction are
used. Three samples across two cells were reissued after upstream
transport errors that produced no model response. Each reissue used
the cell's declared configuration and replaced the failed attempt at
the same sample index; these are infrastructure retries rather than
response-dependent model retries. One further cell was reissued in
whole: the control condition on Claude Sonnet 5 returned an empty
response at one sample index, and all ten samples of that cell were
redrawn under the same declared configuration. Both draws are retained
in the released records and the analysis uses the later one.

For each fixed model--input condition we draw $R=10$ single-pass
samples. Temperature $1.0$ is requested on every execution;
no provider echoes sampling parameters in its response, so the record
is of the requested configuration rather than a provider
confirmation. Locally sampled models additionally use
$\mathrm{top}_p=1.0$ with a per-sample seed recorded; other sampling
parameters for hosted models remain at provider defaults. Outputs containing no action
word never satisfy reproducible correctness; the one such response,
in the reissued cell noted above, is not among the executions
analysed. The
criterion applies only to the declared executions and does not imply
zero probability of another output outside them.

Let $y_r$ denote the operational outcome of execution $r$ and $y^*$
the predefined correct action. Reproducible correctness requires
\begin{equation}
    y_r = y^* \quad \forall r \in \{1,\ldots,R\}.
    \label{eq:reproducible-correctness}
\end{equation}

This yields three operational states. A condition is
\emph{reproducibly correct} when all $R$ executions produce the
intended action, \emph{reproducibly incorrect} when all $R$ executions
produce the same incorrect action, and \emph{non-reproducible} when
the operational decision varies across executions.

Because the prompt requests a single word, most replies are one word
and the operational outcome is that word. Where a model returns prose
instead, we take the action it commits to: the leading action word
when the reply opens with one, and otherwise the final action word,
since such replies restate the question before answering it. Five
model--condition cells contain replies whose leading and final action
words differ; in two of the five the resulting operational state
differs between the two extraction positions. In every other cell both
positions yield the same action in every reply.

\begin{table}[h]
\centering
\caption{Classification of outcome patterns for the car-wash task
($R=10$; correct action: \texttt{DRIVE}).}
\label{tab:operational-states}
\vspace{6pt}
\small
\begin{tabular*}{0.75\columnwidth}{@{\extracolsep{\fill}}ll@{}}
\toprule
\textbf{Observed outcomes} & \textbf{Operational state} \\
\midrule
10 \texttt{DRIVE} / 0 \texttt{WALK} & Reproducibly correct \\
10 \texttt{WALK} / 0 \texttt{DRIVE} & Reproducibly incorrect \\
7 \texttt{WALK} / 3 \texttt{DRIVE} & Non-reproducible \\
\bottomrule
\end{tabular*}
\end{table}

Reliability is evaluated for each fixed model--input condition rather
than assigned to the model in isolation. Aggregate accuracy can obscure
whether a model is consistently correct, consistently wrong, or
variable under one unchanged input; these are operationally different
states.

Supplying task-relevant semantics is the intervention under study, not
a requirement that the model recover omitted semantics unaided. We
refer to constructing and testing natural-language task specifications
against~\eqref{eq:reproducible-correctness} as \emph{Substrate
Engineering}, and to the specification itself as the \emph{substrate}.
A substrate is not a prompt that scores well. It is written, hashed and
versioned, and it is tested against~\eqref{eq:reproducible-correctness}
for each model--condition pair it is claimed to control: for that pair it
either establishes reproducible correctness or is recorded as failing to.


\subsection{Interventions and Measurement}
\label{sec:interventions-measurement}

We treat the input text as the experimental variable. All
interventions modify only the model input: conventional prompts request
reasoning procedures, expertise, or emphasis, whereas the substrate
specifies relations among the task objective and action semantics. A
tenth condition serves as a control: it appends the single sentence
\emph{The correct answer is drive.} to the baseline question. It is not
a prompting technique but a test of whether the decision is reachable
by assertion, and it is reported alongside the others rather than
treated as an upper bound, because it does not establish reproducible
correctness on every model.
Condition text is hashed and recorded with each execution, and the
three reissued samples inherit their cell's declared text and hash;
exact texts, model identifiers, access levels, and inference
configurations are given in the appendix.

The emitted action is the sole basis for assigning reproducible
correctness. Hosted models are evaluated behaviourally only. For
open-weight models, we additionally measure next-token preference and
internal activations using the full output distribution.

Internal quantities are read at the \emph{answer position}: the final
token of the prompt once the model's chat template has been applied,
whose next-token distribution is the model's first answer word. These
measurements are teacher-forced --- a single forward pass over the
prompt, with no generation --- so what is measured is the distribution
the model would sample from rather than a sampled outcome. At that
position we define
\begin{equation}
    M(x) = \log \sum_{t \in \mathcal{D}} P(t \mid x)
         - \log \sum_{t \in \mathcal{W}} P(t \mid x),
    \label{eq:decision-margin}
\end{equation}
where $\mathcal{D}$ and $\mathcal{W}$ contain the measured token forms
of each action, including case and leading-space variants. Positive
$M$ favours \texttt{DRIVE}. This is the logit-difference readout
standard in circuit analysis \citep{wang2023interpretability}, summed
over surface forms because the casing models emit varies. The effect of intervention $u$ relative
to the baseline input is
\begin{equation}
    \Delta M(u) = M(u) - M(\emptyset).
    \label{eq:intervention-effect}
\end{equation}

Activation differences are treated as observational unless directly
intervened on. We attribute a causal effect on the measured preference
only when replacing a baseline internal state changes the downstream
margin. $M$ is supplementary: operational-state labels depend only on
the $R$ emitted actions.


\subsection{Two-Stage Design}
\label{sec:two-stage-design}

The primary behavioural analysis is a transition study over 28 models
that fail the baseline task: under the baseline condition each is
either reproducibly incorrect, naming the wrong action on all ten
samples, or non-reproducible, naming different actions across them.
Cohort membership is fixed from the baseline condition before any
intervention is evaluated.

Stage one evaluates ten conditions through repeated single-pass
executions: the baseline task, seven conventional prompt
interventions, one shared substrate, and a control that states the
correct action in the prompt. Each model--input condition is classified by the criterion
in Section~\ref{sec:operational-reproducibility} before any internal
analysis is performed.

Stage two studies the baseline-to-substrate transition in open-weight
models by \emph{activation patching}, the causal-mediation technique
established for language models by \citet{vig2020causal} and
\citet{meng2022locating} and applied to circuit analysis by
\citet{wang2023interpretability}. We run the substrate-conditioned
prompt and record the residual-stream state at the answer position at
every layer. We then run the baseline prompt and, at one layer, replace
the residual at that same position with the recorded state, leaving
every other position and every other layer untouched and completing the
baseline computation. No substrate text is present in the patched
context. We read the final decision margin, and repeat the procedure at
each layer.

Two aspects differ from the most common form of the technique. The
source and the receiver are two natural prompt conditions rather than a
clean run and a corrupted one, so no noised or ablated input enters the
comparison \citep{zhang2024towards}. And the intervention is anchored to
a behavioural transition established in advance: internals are
interpreted only for model--condition pairs whose operational state is
already fixed at both ends by
Equation~\eqref{eq:reproducible-correctness}, rather than for a single
successful generation.

A change in margin establishes a causal effect of the transplanted
state on the measured preference; a sign reversal establishes that the
state is sufficient, under the remaining baseline computation, to
reverse that preference. The intervention is restricted to one token
position, rather than swept over positions as in the canonical
causal-tracing map \citep{meng2022locating}: a state that mattered only
through the question span, or through several positions jointly, would
not be detected by it. It does not establish reproducibly correct
behaviour, localise where the decision is formed
\citep{zhang2024towards}, or identify the effect of an individual
constraint.
